\documentclass[12pt]{article}

% House style preamble
\input{.house-style/preamble.tex}

% Conceptual metaphor formulas (M&S convention: quoted uppercase italics)
\newcommand{\metaphor}[1]{\enquote{\textit{\MakeUppercase{#1}}}}

% Update PDF metadata
\hypersetup{
  pdftitle={The Projectibility of Metaphor: Mechanism Support and the Reliability of Cross-Domain Inference}
}

\title{The Projectibility of Metaphor:\\Mechanism Support and the Reliability of Cross-Domain Inference}
\author{Brett Reynolds \orcidlink{0000-0003-0073-7195}%
\thanks{Contact: \href{mailto:brett.reynolds@humber.ca}{brett.reynolds@humber.ca}}\\
Humber Polytechnic \& University of Toronto}
\date{}

\begin{document}
\maketitle

\begin{abstract}
What separates a productive metaphor from a superficial one is the projectibility of its cross-domain inferences: what observing some features of a mapping licenses us to predict about the rest \citep{Goodman1955}. This paper argues that metaphorical projectibility is underwritten by the worldly relations that hold property clusters together in both source and target domains, and that these come in grades of support that \posscite{lakoff1980} Conceptual Metaphor Theory leaves untheorized. \posscite{boyd1993metaphor} Homeostatic Property Cluster (HPC) theory names the strongest grade, where corrective feedback maintains the cluster, but it is one grade among several. Recent extensions of CMT bear on adjacent questions~-- which metaphor context selects (Kövecses), whether it's used deliberately (Steen)~-- but leave projectibility untouched. On the view developed here, metaphor is the cross-domain extension of a projectibility profile on partial property overlap, and a mapping projects when the shared relational structure is supported on both sides. This recasts several standing disputes in metaphor theory and predicts the full range from literary metaphor through everyday to theory-constitutive, on a single dimension: the degree to which mechanisms in both domains sustain projectible inference. A case study of canine trail cognition illustrates the framework and generates a testable prediction: if \term{trail} is a projectibility profile, dogs should recruit it for objects sharing only visual-geometric trail properties, even across functional domains.
\end{abstract}

\noindent\textbf{Keywords:} conceptual metaphor theory, projectibility, projectibility profile, support grades, homeostatic property cluster, theory-constitutive metaphor, cross-domain extension, comparative cognition

\section{Introduction}
\label{sec:intro}

My dog follows painted lines on soccer fields.%
\footnote{An observation, not an experiment. Section~\ref{sec:dogs} discusses confounds.} Not occasionally or ambiguously: he tracks along the white line the way he tracks along a trail through the woods. There's no scent concentration on painted lines. There's no substrate difference (paint on grass), no energy advantage, no human to follow. The painted line shares only the visual-geometric properties of a trail: a linear, bounded, high-contrast strip extending through space. Everything else that makes a trail a trail is absent. It is, in short, a trail with all the interesting parts politely removed.

What's he up to? The obvious answer is \enquote{following a trail.} But a painted line isn't a trail. It belongs to a different functional domain (spatial demarcation, not locomotion). If my dog is treating it as a trail, he's extending a category across domains on the basis of partial property overlap. That's a big \enquote{if,} and I don't have controlled evidence to resolve it. But the observation is puzzling enough to motivate a question: what kind of theoretical framework would make sense of cross-domain category extension, in dogs or in humans?

That question arrived from one direction. From another direction entirely, I'd been thinking about how metaphor works in theory construction (prompted by \posscite{mitchell2024} analysis of AI metaphors, discussed on the \emph{Many Minds} podcast; \citealt{cooperrider2026}) and realized that conceptual metaphors have the structure of homeostatic property clusters. The two threads converged on the same framework. That convergence is what motivates the paper: the theoretical synthesis wasn't reverse-engineered from a single dog observation, and the dog observation wasn't generated by the theory. They met in the middle. My dog, of course, had no interest in the synthesis.

The framework draws on two of the most influential accounts of non-classical categorization, developed in parallel without mutual engagement. Both pulled metaphor out of the literary tradition: Lakoff and Johnson into everyday cognition, Boyd into scientific theory construction. One consequence of the profile reframing is that it brings literary metaphor back into the picture (Section~\ref{sec:literary}).

\textcite{lakoff1980} proposed that abstract thought is fundamentally metaphorical, grounded in embodied experience: we understand \term{argument} in terms of \term{war}, \term{time} in terms of \term{space}, \term{love} in terms of \term{journey}. \textcite{lakoff1987} extended this into a comprehensive theory of categorization built on image schemas, radial categories, and idealized cognitive models. Independently, \textcite{boyd1979} argued that theory-constitutive metaphors play an irreplaceable role in scientific theory construction by providing epistemic access to natural kinds whose properties aren't yet fully understood. Boyd's Homeostatic Property Cluster (HPC) theory \citep{boyd1989, boyd1999} gave this a general framework: natural kinds are individuated not by essences but by clusters of properties maintained by causal mechanisms, and the clustering is \enquote{important for scientific explanation or for the formulation of successful inductive inferences} \citep[526]{boyd1993metaphor}.

\textcite{lakoff1987} doesn't cite Boyd. Boyd doesn't cite Lakoff. The experimental metaphor literature \citep{flusberg2018, flusberg2024, thibodeau2019, holyoak2018} doesn't cite Boyd either. HPC thinking has reached linguistics through other channels: \textcite{croft2000} applies evolutionary thinking to linguistic categories in ways that share HPC's rejection of classical definitions, and the mechanistic property cluster (MPC) tradition in psychiatry \citep{kendler2011} provides precedent for multi-level mechanisms in non-biological domains. Hofstadter and Sander's (\citeyear{hofstadter2013}) thesis that analogy is the core of all cognition, including categorization, is compatible but addresses a different question: what the cognitive operation is, not what makes some extensions persist while others die. The specific synthesis proposed here (a projectibility-profile reframing of conceptual metaphor theory, with graded support underwriting projectibility and conventionalization) has not been attempted. The convergence between Lakoff and Boyd is striking: both frameworks reject classical categories, both invoke property clustering, both treat graded typicality as genuine structure rather than noise. But they locate the explanatory work in different places. Lakoff explains how metaphors arise (embodied schemas); Boyd explains why they succeed (accommodation to causal structure). Neither provides the full picture.

The reframing supplies the stability that Lakoff and Johnson's account needs but doesn't provide, by treating a source domain as a projectibility profile: a clustered pattern of properties whose partial observation licenses inference, held together by support of some grade. On this view, metaphor is the cross-domain extension of such a profile on partial property overlap. A source-domain category is extended to instances in a different functional domain that share only some of its properties. The extension projects when the shared relational structure is supported in the target domain, not just in the source. It fails when the overlap is superficial and nothing in the target sustains it. Boyd's Homeostatic Property Cluster (HPC) theory names the strongest grade of that support, where corrective feedback maintains the cluster; but homeostatic control is one grade among several \citep{reynolds2026kindsProjectibilityProfiles}, and most metaphors that work do so on weaker grades.

The paper's contribution is an account of metaphorical projectibility: what makes some cross-domain inferences reliable and others not. The dog observation motivates and illustrates. First, I show that Lakoff and Johnson's source domains are property clusters without an adequate story about what maintains them (Section~\ref{sec:source-domains}), and that recent extensions of CMT bear on adjacent questions but leave projectibility untheorized (Section~\ref{sec:situating}). I then show that Boyd's apparatus supplies the missing story (Section~\ref{sec:hpc-tcm}), graded into support that runs from bare stabilization to strict homeostatic control, and that combining them reframes metaphor as cross-domain profile extension (Section~\ref{sec:reframing}). Section~\ref{sec:consequences} draws out the consequences: what each framework gains, which disputes it recasts, and what the reframing predicts. Section~\ref{sec:literary} shows that the framework handles the full range of metaphor, from literary through everyday to constitutive, predicted by a single dimension: support grade. Then I return to the dog (Section~\ref{sec:dogs}). Canine trail cognition turns out to be richly documented: dogs possess a functional spatial category \term{trail}, a genuine projectibility profile maintained at the controller grade by ecological feedback. This establishes that non-human categories can have profile structure. The framework then generates a prediction: if metaphor is cross-domain profile extension, dogs should recruit \term{trail} for objects sharing only visual-geometric trail properties. That prediction is testable but untested. Section~\ref{sec:phylogeny} considers what's distinctively human about the superstructure built on this basic operation.

\section{Source domains without mechanisms}
\label{sec:source-domains}

Lakoff and Johnson's Conceptual Metaphor Theory begins with a striking observation about ordinary language. We don't just talk about arguments in terms of war~-- we experience them that way. \enquote{Your claims are indefensible,} \enquote{He attacked every weak point,} \enquote{I demolished his argument} aren't decorative; they structure how we think about argumentation \citep[4]{lakoff1980}. The metaphorical mapping is systematic: positions in arguments correspond to positions on a battlefield; strategies correspond to tactics; winning corresponds to defeating the opponent.

The theoretical apparatus behind this observation is built on property clusters, though Lakoff and Johnson don't use that term. Their \enquote{experiential gestalts} are clusters of co-occurring properties that emerge from recurring bodily experience. Causation, for example, involves a cluster of twelve properties including agency, volition, physical contact, change of state, and temporal priority \citep[69--72]{lakoff1980}. No single property is necessary or sufficient. Members of the category \term{causation} share overlapping subsets, producing the graded membership and prototype effects that \textcite{rosch1975} documented.

\textcite{lakoff1987} extends this into a full theory of categorization. Image schemas (\term{container}, \term{path}, \term{link}, \term{part-whole}, \term{centre-periphery}) are \enquote{relatively simple structures that constantly recur in our everyday bodily experience} \citep[267]{lakoff1987}. They're directly meaningful because they're \enquote{directly and repeatedly experienced because of the nature of the body and its mode of functioning in our environment} \citep[267]{lakoff1987}. The \term{source-path-goal} schema, for instance, captures the structure of every movement from one place to another: a starting point, a trajectory, a destination, a direction. This schema then grounds metaphors like \metaphor{purposes are destinations} (\enquote{she's going somewhere in life,} \enquote{we've gotten sidetracked}).

Radial categories elaborate the picture further. Lakoff's analysis of Dyirbal noun classification shows categories linked by chains of motivated extensions: women, fire, dangerous things, and sundry other items clustered by domain-of-experience principles, myth-and-belief links, and image-schema transformations \citep[Ch.~6]{lakoff1987}. Central members have the most properties; peripheral members fewer. The structure resists classical definition. \enquote{The general principles given make sense of Dyirbal classification, but they do not predict exactly what the categories will be} \citep[Ch.~6]{lakoff1987}.

All of this is largely right. Lakoff was right that categories aren't classical. He was right that metaphor does cognitive work, not just decorative work. He was right that embodiment matters. The embodied simulation tradition that grew from this work offers a genuine mechanism story: shared bodies produce shared sensorimotor simulations, which ground shared conceptual structure. Because all humans walk upright, all humans have a \term{source-path-goal} schema; because all humans exert force, all humans have a \term{force dynamics} schema. The shared embodiment explains why metaphorical mappings converge across individuals.

The problem is what this story leaves unexplained. Shared bodies explain why \emph{individuals} converge on the same schemas. They don't explain why \emph{communities} maintain specific metaphorical mappings over generations, why different communities with the same bodies metaphorize differently, or why mappings erode in structured ways when cultural transmission weakens (as in Dyirbal, below). Bodies haven't changed in the communities where Dyirbal is dying; what's weakened is the cultural mechanism transmitting the category system. CMT has plausible candidate mechanisms at the individual level (neural recruitment, embodied simulation, associative learning). What it lacks is a principled account tying community-level persistence, cross-generational transmission, and the reliability of projected inferences into a single explanatory structure. The mechanisms Lakoff needs are ecological and social, not just neural.

The gap shows up most clearly in stability. What maintains radial category structure over time, across speakers, across generations? Lakoff describes the chaining principles (domain-of-experience, metonymy, image-schema transformation) but doesn't theorize the mechanisms that keep the structure in place. Image schemas are said to emerge from bodily experience, but there's no account of why they're stable across individuals or what happens when they fail.

Language death provides the most vivid evidence for the gap. \posscite{schmidt1985} study of Dyirbal in decline shows peripheral links breaking off first, leaving central members \citep[as discussed in][Ch.~6]{lakoff1987}. Radial structure collapses from edges inward. This is gradient stability loss, and it cries out for a mechanistic explanation. The same pattern appears in English countability: as the singulative \mention{datum} disappears, tight properties of \mention{data} (cardinal selection, \mention{a(n)}) erode before loose ones (agreement, \mention{many}), producing asymmetric fraying rather than uniform collapse \citep[Ch.~9]{reynolds2025hpcbook}. Lakoff notes the Dyirbal pattern but has no theory of the mechanisms whose weakening produces it.

\textcite{murphy1996} identifies a related problem. If source domains literally structure target concepts (strong CMT), then multiple metaphors for the same concept create incoherence: \enquote{it is difficult to see how these two metaphors could simultaneously inform the same concept} when \term{love} is structured by both \term{journey} and \term{war} \citep[186]{murphy1996}. Murphy's solution is that metaphor works through structural similarity between independently represented concepts. Each metaphor \enquote{simply picks out different aspects of the concept's content} \citep[196]{murphy1996}. This is the right move, but it needs a story about what makes the structural similarities real rather than arbitrary.

The 2003 Afterword to \emph{Metaphors We Live By} gestures toward mechanisms. Lakoff and Johnson invoke neural recruitment (\enquote{Neurons that fire together wire together,} 2003:~256) and stable conventional systems. But neural recruitment, while it does maintain associations at the individual level (long-term potentiation sustains co-activation patterns), doesn't explain community-level stability or cross-generational transmission. A Hebbian account tells us why \emph{I} maintain \metaphor{argument is war}; it doesn't tell us why English speakers as a community maintain it, why it persists across generations, or why some metaphorical mappings conventionalize while others die. Lakoff has phenomenology without etiology.

% === BEGIN inserted Section: Situating the account (draft, 2026-06-24, projectibility-first). Source-grounded; prose still needs Brett's voice pass. Mirror of notes/section3-situating-DRAFT.tex ===
\section{Situating the account}
\label{sec:situating}

Conceptual Metaphor Theory has not stood still since 1980, and two developments speak directly to the gap I've described. Neither closes it, but saying why locates the contribution. The question this account puts at the centre is the projectibility of metaphorical inference: what observing some features of a mapping licenses us to predict about the rest \citep{Goodman1955}. I take \term{projectibility} in the profile sense of \textcite{reynolds2026kindsProjectibilityProfiles}: a category earns standing by tracking a worldly pattern whose partial observation licenses inferences under specified conditions. Which cross-domain mappings support reliable inference, and what underwrites that reliability? Extended conceptual metaphor theory and deliberate metaphor theory answer different questions, and the three can be held apart.

\subsection{Extended conceptual metaphor theory}

\textcite{kovecses2017, kovecses2020} extends CMT along two axes. First, a single conceptual metaphor is reconstrued as a hierarchy across four levels of schematicity~-- image schema, domain, frame, and mental space~-- rather than a mapping between two domains \citep{kovecses2017}. This dissolves a long-standing grain problem: source domains carry more material than any mapping uses, and locating the metaphor on the hierarchy says which material is in play. Second, metaphor in discourse is given a contextual component. Kövecses distinguishes four context types (situational, discourse, bodily, and conceptual-cognitive), each with factors that \enquote{prime} a speaker toward one metaphor rather than another in a given situation \citep{kovecses2015, kovecses2020}. This answers a real complaint: classical CMT, built from dictionaries and decontextualized examples, had little to say about why a particular metaphor surfaces in particular discourse.

Both moves operate within the synchronic life of a metaphor, and neither asks whether the mapping supports projectible inference. The contextual component is a theory of production-side selection: which stored mapping a speaker reaches for, here and now, under contextual pressure. And the dynamism Kövecses adds is the dynamism of processing, not of history. The offline/online distinction in extended CMT is a distinction between memory systems: image schemas, domains, and frames are \enquote{conventionalized knowledge structures in long-term memory,} while mental spaces are \enquote{online representations of our understanding of experience in working memory} \citep[326]{kovecses2017}. Conventionalization enters as a property a structure already has (stored, shared across a speech community), not as something the theory explains, and reliability of inference doesn't enter at all. Extended CMT takes the supraindividual layer~-- the decontextualized patterns a language and culture happen to share~-- as given, and theorizes what speakers do with it online. Which of those mappings project, and what keeps them projectible across generations, is a different question. Mechanism support answers it: a mapping projects when the shared relational structure is causally maintained in both domains, not just in the source. The two layers nest rather than compete. Contextual priming selects among mappings that are already available; projectibility, underwritten by mechanisms, is what made them worth keeping.

The difference is testable, not terminological. Extended CMT doesn't predict that peripheral, weakly maintained properties erode before core ones when transmission weakens (the Dyirbal and \mention{data} cases, Section~\ref{sec:source-domains}); it doesn't predict that a high-frequency mapping fails to conventionalize when its target domain has no mechanisms sustaining the shared structure (Prediction~1); and it makes no prediction about non-linguistic animals (Section~\ref{sec:dogs}). The projectibility account does. And the gap is one Kövecses partly concedes from the inside: his processing model is, in his own words, \enquote{informal and hypothetical,} and \enquote{it is only through psychological experiments that the validity of this \ldots\ model could be confirmed} \citep[125]{kovecses2020article}. Grounding reliability of inference in world-side mechanisms is one way to give the framework external traction it doesn't yet have.

\subsection{Deliberate metaphor}

\textcite{steen2008} adds a third dimension to the linguistic and conceptual ones: a metaphor can be used \emph{deliberately}, as a metaphor that points the addressee to a source domain, or non-deliberately, as part of the ordinary structure of the language, where \metaphor{argument is war} passes unnoticed. Deliberateness is a fact about how a metaphor functions in communication. It cross-cuts projectibility rather than tracking it, and the two axes are worth separating explicitly.

\begin{table}[ht]
\centering\small
\begin{tabular}{@{}>{\raggedright\arraybackslash}p{.20\textwidth}>{\raggedright\arraybackslash}p{.36\textwidth}>{\raggedright\arraybackslash}p{.36\textwidth}@{}}
\toprule
 & Deliberate & Non-deliberate \\
\midrule
High projectibility & A theory-constitutive metaphor a scientist knowingly deploys (\metaphor{the brain is a computer}) to direct research & A constitutive mapping so entrenched it's gone invisible in the field that lives by it \\[6pt]
Low projectibility & A fresh literary metaphor (Dickinson's \enquote{Hope is the thing with feathers}) & A dead idiom whose source no longer projects anything \\
\bottomrule
\end{tabular}
\caption{Deliberateness (Steen) and projectibility cross-cut. Mechanism support underwrites the rows, not the columns.}
\label{tab:deliberate-projectibility}
\end{table}

The three-way typology developed below (literary, everyday, constitutive; Section~\ref{sec:literary}) is the projectibility axis of Table~\ref{tab:deliberate-projectibility}; deliberateness is the other. Because the axes are independent, whatever processing signature deliberateness predicts, projectibility predicts a separate one: how reliably the mapping supports inference, how gracefully it degrades when source-domain knowledge is disrupted, and whether it conventionalizes over time. The two can be measured apart. That extended CMT itself lists incorporating deliberate metaphor among its open problems \citep[113]{kovecses2020article}, rather than deriving it from its apparatus, is a sign the communicative axis really is separate from the conceptual one. I take the distinction DMT draws, not its further claim that deliberate metaphors capture more attention and are better remembered~-- an entailment that \textcite{gibbs2015dmt} and \textcite{thibodeau2017} find poorly supported, as \textcite{flusberg2018} note.

\subsection{Cultural and cross-linguistic variation}

Communities with the same bodies metaphorize differently, and the differences are patterned, not free (Section~\ref{sec:source-domains}). Embodiment alone struggles here: if shared sensorimotor experience grounds the mappings, shared bodies should converge. The projectibility account locates the variation in the mechanisms rather than the bodies. Where ecologies, institutions, and practices differ across communities, the property clusters they maintain differ too, and so do the mappings that stay projectible. This gives a mechanism reading of the cultural variation \textcite{kovecses2005} documents: variation tracks differences in what the world, locally, keeps clustered, which also predicts which parts of a mapping should project across communities and which should fray. The point connects to discourse-based and cross-linguistic metaphor analysis \citep[e.g.][]{semino2008, musolff2006}, which shows metaphors doing situated cultural work that a purely embodied account underdetermines. Projectibility says which of those situated mappings have the causal backing to support reliable inference and persist.

\subsection{Concept structure and analogy}

The account already treats structure-mapping systematicity \citep{gentner1983} and the purpose constraint \citep{holyoak2018} as components of cross-domain extension (Section~\ref{sec:consequences}). It connects to an older claim about category structure. \textcite{murphymedin1985} argued that what makes a concept cohere isn't similarity among its features but the background theories a learner brings: \enquote{concepts are coherent to the extent that they fit people's background knowledge or naive theories about the world,} theories that \enquote{help to relate the concepts in a domain and to structure the attributes that are internal to a concept} \citep[289]{murphymedin1985}. Similarity, they argue, is too unconstrained to pick out which concepts are coherent~-- which is, in Goodman's terms, the projectibility problem for concepts. The profile account is the world-side completion of that move. A learner's theory projects to new instances to the degree it tracks causal mechanisms in the world that maintain the cluster. This separates the account from similarity-based models (prototype and exemplar), where cohesion lives in the representation, and from structure-mapping read as purely formal alignment: what makes an alignment support reliable inference is causal maintenance in both domains, not the match itself.
% === END inserted Section: Situating the account ===

\section{Projectibility, support, and theory-constitutive metaphor}
\label{sec:hpc-tcm}

\textcite{boyd1979} developed his account of theory-constitutive metaphors in the context of scientific realism, not cognitive science. But the apparatus he built turns out to be exactly what metaphor theory needs.

A theory-constitutive metaphor (TCM) is a metaphorical expression that constitutes \enquote{an irreplaceable part of the linguistic machinery of a scientific theory} for which \enquote{no adequate literal paraphrase is known} \citep[360]{boyd1979}. The canonical examples come from cognitive psychology: thought as \enquote{information processing,} brain as \enquote{computer,} memory as having \enquote{storage,} \enquote{indexing,} and \enquote{retrieval} \citep[486]{boyd1993metaphor}. These aren't pedagogical metaphors that explain established theories; they're constitutive metaphors that make certain questions askable in the first place.

The key mechanism is accommodation to causal structure. Boyd's central claim is that metaphor succeeds by \enquote{introducing terminology, and modifying usage of existing terminology, so that linguistic categories are available which describe the causally and explanatorily significant features of the world} \citep[358]{boyd1979}. Metaphors don't just illustrate; they refer. And referential success depends not on explicit definitions but on epistemic access: the capacity of a linguistic community to detect and discover facts about the referent through ongoing use of the metaphorical term.

This is where HPC theory enters. \textcite[526--527]{boyd1993metaphor} characterizes homeostatic property cluster kinds through a set of features, of which five are central to the present argument. First, there's a family of properties that contingently cluster in nature. Second, their co-occurrence results from \enquote{what may be metaphorically (sometimes literally) described as a sort of \emph{homeostasis}} (mechanisms that favour the presence of some properties given others) \citep[526]{boyd1993metaphor}. \emph{Homeostatic} is doing specific work here, and it's worth keeping the grades of stabilization apart, since the term is easily stretched to cover achievements that differ \citep{reynolds2026kindsProjectibilityProfiles, reynolds2026notEveryStableCluster}. A property cluster may merely be \term{stabilized}: its co-availability persists across a range, as a valley's shape persists under erosion and gravity, with no active process tracking departures. It may be actively \term{maintained}: a process preserves or regenerates it, as healing closes a wound. Or it may be \term{controlled} in the strict sense, by perturbation-sensitive feedback that restores a higher-scale relation through lower-scale variation, as corrective loops hold body temperature. \term{Homeostatic} belongs to this last, controller grade: feedback where the presence of some properties promotes the presence of others, making the cluster self-reinforcing and robust to perturbation. The grades matter for metaphor, because different mappings are underwritten at different grades and only some reach the controller grade. What they share, and what does the work throughout, is \term{support}: the worldly relations that make partial observation of a cluster license inference about the rest. Throughout this paper, \term{mechanism} refers to these world-side causal processes (foot traffic, scent deposition, institutional norms), not to cognitive mechanisms in the organism; where cognitive processing is relevant, I say so explicitly. Third, the clustering is causally important: significant effects are produced by the conjoint occurrence of properties together with the underlying mechanisms. Fourth, the kind is \enquote{important for scientific explanation or for the formulation of successful inductive inferences} \citep[526]{boyd1993metaphor}~-- the property I'll call \term{projectibility}: membership in the kind supports reliable predictions about unexamined instances. Fifth, no refinement replacing the kind term with a more extensionally precise term \enquote{will preserve the naturalness of the kind referred to} \citep[527]{boyd1993metaphor}. Apparent vagueness isn't imprecision: \enquote{vagueness in extension is by no means always indicative of any imprecision at all. For some kinds, `vagueness' in the application of the associated terminology is precisely indicative of \emph{precision} in the accommodation of language to the actual structure of the world} \citep[530]{boyd1993metaphor}.

(A clarification: \textcite[526]{boyd1993metaphor} takes extensional vagueness to be irresolvable even given all relevant facts. I think the boundaries are sharp but epistemically inaccessible, the gradience reflecting our measurement rather than the category itself \citep{reynolds2025sorites, reynolds2025hpcbook}. This divergence doesn't affect the argument that follows, which depends on property clustering and mechanisms, not on the metaphysics of boundary vagueness.)

TCMs display \enquote{inductive open-endedness} \citep[363--364]{boyd1979}. They're \enquote{invitations to future research} whose function is to put researchers \enquote{on the track of these respects of similarity or analogy; indeed, the employment of terms in such metaphors may best be understood as contributing to features of the world delineated in terms of those~-- perhaps as yet undiscovered~-- similarities and analogies} \citep[489]{boyd1993metaphor}. The inability to fully explicate a TCM isn't a defect. When the metaphor refers to an HPC phenomenon, complete explication may be impossible because \enquote{the properties that constitute the homeostatic property cluster may not even be finite in number and they may vary significantly from time to time or place to place} \citep[529]{boyd1993metaphor}.

This connects to catachresis~-- Black's term for the use of metaphor to fill vocabulary gaps \citep[356--357]{boyd1979}. TCMs are catachrestic not because they're substitutes for literal language but because they have \enquote{distinctive capacities and merits} that literal language can't replicate \citep[482]{boyd1993metaphor}. The term \enquote{homeostatic property cluster} when applied to grammatical categories is catachrestic in exactly this sense. There was no literal expression for \enquote{a category constituted by co-occurring properties sustained by causal mechanisms whose boundaries resist explicit specification.} \enquote{Prototype} came closest, but it smuggles in Lakoff's cognitive framing. TCMs are \enquote{conceits}~-- they extend \enquote{not through one literary work, but through the work of a generation or more of scientists} and become \enquote{the property of the entire scientific community} \citep[361]{boyd1979}.

Applying HPC theory to grammar \emph{is} a theory-constitutive metaphor in Boyd's sense. It imports inferential machinery (homeostasis, property co-occurrence, maintained boundaries, causal mechanisms) from biology into linguistics. The metaphor does constitutive work: it makes certain questions askable~-- What are the homeostatic mechanisms? Why do these properties cluster?~-- that weren't askable under classical-category or prototype-category frameworks. The partiality of the mapping is where the theoretical work happens.

The reflexivity also disciplines the framework. If HPC-applied-to-metaphor is itself a TCM, it has to meet the success conditions it proposes: identifiable mechanisms, projectibility to new cases, exclusion of non-cases. It would fail if no identifiable homeostatic mechanisms maintained the property clusters in either source or target domains (rendering \enquote{HPC kind} vacuous), if the framework couldn't discriminate productive from superficial metaphors beyond what existing theories already capture, or if every category extension counted as cross-domain (draining \enquote{functional domain} of content). These are genuine failure conditions, not escape clauses.

\section{Metaphor as cross-domain profile extension}
\label{sec:reframing}

Metaphor is the extension of a projectibility profile to instances in a different functional domain where only some of the profile's clustered properties are present. The source domain is such a profile: a family of properties whose clustering supports inductive inference, held together by support of some grade (from bare stabilization through active maintenance to homeostatic control). The target domain shares some of those properties (enough to trigger recruitment of the source-domain category) but lacks the full support suite. Metaphorical inference is reliable when the shared relational structure is supported in the target domain, not just in the source. It is unreliable when the overlap is superficial and nothing in the target sustains it.

The notion of \enquote{functional domain} needs to be made precise, since the entire literal/metaphorical distinction depends on it. Two instances belong to the same functional domain when the same suite of supporting mechanisms maintains their property clusters. The test is which mechanisms maintain the trail-relevant properties (visual contrast, substrate, scent, directional use), not which mechanisms explain the object's existence.

A forest trail and a busy sidewalk are in the same domain. Yes, someone poured the concrete (planning), but the trail-relevant properties are maintained by locomotion: heavy use wears the surface, deposits scent, creates visual contrast at the edges, and sustains directional traffic. A lightly used sidewalk is a worse trail (less scent, less wear, less traffic) for the same reason an overgrown forest path is: the locomotion mechanism is weak. Desire paths make the point decisively. When pedestrians ignore the planned sidewalk and cut across grass, a trail forms through exactly the mechanisms that maintain forest trails (foot traffic → vegetation wear → visual contrast → more traffic). The planning mechanism tried to put the trail elsewhere; the locomotion mechanism built one anyway. Anyone who has watched a campus lawn lose to students has seen the same mechanism in miniature. That's two mechanisms in opposition, producing two different trails. The functional domain is individuated by which mechanism maintains the property cluster, and the desire path shows that locomotion and planning are genuinely distinct mechanism suites.

A painted line on a soccer field is in a different domain. No amount of walking on it makes it more line-like. Its visual contrast is maintained by paint application and rule enforcement, not by use. Its linear extent reflects spatial-demarcation planning, not locomotion. If the paint fades, the line disappears regardless of how much traffic crosses it. None of the trail-maintaining feedback loops operate. The property profiles partially overlap (linear extent, visual contrast, bounded width), but the maintaining mechanisms don't.

Cross-domain extension occurs when a projectibility profile recruits an instance whose properties are maintained by a different mechanism suite. This criterion is observer-independent in the relevant sense: what mechanisms maintain a property cluster is a fact about the world, not a judgment call by the theorist, even though identifying those mechanisms requires empirical investigation.

Take the metaphor \metaphor{argument is war}. Both argument-kinds and war-kinds exhibit competition for dominance, strategic reasoning, alliance formation, victory and defeat. These are shared relational properties, not shared mechanisms. The mechanisms differ: the property cluster in war is maintained by military institutions, training, and logistics; in argumentation, by discursive norms, social hierarchies, and reputational stakes. What makes the metaphor productive is that the shared relational structure (competition, strategy, alliance) is maintained in both domains, each by its own mechanism suite. Inferences licensed by the relational structure in one domain (flanking manoeuvres, fortified positions, tactical retreat) are projectible to the other. The metaphor is stable because the mechanisms maintaining each cluster persist across communities.

Compare \metaphor{society is an organism}. The structural overlap is genuine and was historically productive: differentiated parts serving specialized functions, regulation of internal states, growth, vulnerability to pathology. Spencer and Durkheim built research programmes on it. But the mechanisms maintaining biological homeostasis (immune response, metabolic regulation, cellular repair) don't operate in societies. Social stability is maintained by institutions, norms, and power structures: a different mechanism suite entirely. The parts of the metaphor that projected from biology (functional differentiation, homeostatic regulation) worked precisely because they mapped onto relational structure maintained in both domains. The parts that didn't (disease, senescence, death of the organism) failed because no social mechanisms sustained them. The framework diagnoses the partial success: mechanism support was real for some relational properties and absent for others.

\textcite{cummiskey1992} provides the sharpest illustration of this success/failure distinction. Caloric theory had phenomenal success (heat transfer experiments worked), but the caloric-as-fluid metaphor couldn't be articulated mechanistically. It failed. Kinetic theory had phenomenal success \emph{and} the metaphor was articulable (molecular motion as mechanism). It succeeded. \enquote{When a new explanatory mechanism can explain the successes and can be more fully articulated, it is reasonably concluded that no explanatory significant feature of the world was picked out by the central term of the theory} \citep[35]{cummiskey1992}. Mechanism articulability discriminates productive metaphors from superficial ones.

\section{Consequences and predictions}
\label{sec:consequences}

This explains what each framework was missing.

The profile account adds three things to metaphor theory. First, support. Lakoff and Johnson describe the clusters; the account explains why they cohere and why some properties are more central than others. Metaphors persist because they map between domains whose property clusters are held in place by support of some grade. When the support weakens, the metaphor dies, just as Dyirbal radial categories collapse from edges inward when cultural transmission weakens.

Second, falsification structure. Lakoff's ICMs are notoriously hard to falsify \citep{glynn2010}. There's no principled way to say \enquote{this isn't a radial category} in his framework. The profile account provides exclusion power through failure modes: if you can't identify mechanisms, or if the cluster doesn't project to new cases, you don't have a genuine profile. Applied to metaphor: if you can't identify mechanisms maintaining the shared relational structure in each domain, or if the metaphorical mapping doesn't support reliable inference, the metaphor is superficial.

Third, graded typicality grounded in mechanisms rather than cognition. Prototype effects aren't cognitive illusions or representational artefacts. They reflect real variation in how many of the cluster's properties are instantiated and how far an item sits from the category core. A metaphorical extension where only a few properties overlap is genuinely less typical than a core instance where many overlap~-- not because membership comes in degrees, but because stability does.

Metaphor adds two things to the profile account. First, cross-domain extension. Profile analysis has focused on within-domain categorization: biological species, chemical kinds, perhaps social kinds. Metaphor shows profiles reaching across domains when partial property overlap triggers recruitment. This is a capacity the account hasn't previously been asked to handle. Second, gradience at the literal/metaphorical boundary. The transition from peripheral literal categorization (\enquote{that sidewalk is a trail}) to metaphorical extension (\enquote{that painted line is a trail}) is gradient (fewer properties instantiated, weaker support) and the point at which it crosses from literal to metaphorical can't be specified independently of investigative context. This is a prediction the account makes but hasn't previously delivered on.

The reframing recasts three standing disputes in metaphor theory. The comparison-versus-categorization debate \citep{gentner1983, glucksberg2001} turns out to capture two stages of profile extension. Novel metaphors are processed as comparisons (structural alignment testing whether the shared structure supports projectible inference in the target domain); conventional metaphors are processed as categorizations (membership in a generalized profile). The career of metaphor \citep{bowdle2005} is a stabilization trajectory: novel metaphors whose relational structure proves to have support in the target domain conventionalize; those without it remain novel or die. What the account adds beyond Bowdle and Gentner's own account is a prediction about \emph{which} novel metaphors will conventionalize (see Prediction~1 below).

The divergence from \posscite{gentner1983} systematicity principle is sharpest in diachronic predictions. Systematicity predicts which relational structures will transfer productively at any given moment: those with richer, more interconnected relations. The profile account predicts something systematicity cannot: which categories will erode when support weakens, and in what order. The Dyirbal case (Section~\ref{sec:source-domains}) illustrates this directly: peripheral members with fewer supporting connections lose membership first; core members with many properties and strong support persist longest. Systematicity is a synchronic processing principle; support is a diachronic stability principle. The two come apart whenever a richly systematic relational structure lacks the causal maintenance to persist. The divergence is visible synchronically too. English \mention{furniture} has richly interconnected semantic structure (functional unity, spatial co-occurrence, shared purchase contexts), but no morphosyntactic mechanism enforces countability distinctions \citep[Ch.~9]{reynolds2025hpcbook}; systematicity predicts the relational richness should do work, and it doesn't.

The attribute-versus-relational mapping dispute \citep{murphy1996, gentner1983} is resolved by the distinction between properties and the support that maintains them. What transfers in metaphor is relational structure (supported in both domains), not intrinsic properties. \textcite[163]{gentner1983} \enquote{systematicity principle} (interconnected relational systems are preferentially mapped) is supported coherence by another name: the relational structures that preferentially transfer are exactly those sustained by causal processes in both source and target.

The embodied-versus-abstract dispute (Lakoff versus Murphy/Gentner) is resolved by recognizing that embodied schemas are one mechanism type among several. The account allows multiple stabilisers: perceptual, motor, social, functional, cultural, institutional. Lakoff was right that bodily experience matters. He was wrong that it's the whole story. Boyd was right that causal structure in the world constrains which metaphors succeed. He was wrong to focus exclusively on scientific metaphor. The dog case developed below actually demonstrates the point: \term{trail} is maintained by deeply embodied mechanisms (walking, sniffing, substrate sensing), and the category's structure reflects the dog's bodily engagement with its environment. The account doesn't replace embodiment. It situates embodiment as one mechanism type among the several that maintain property clusters.

\posscite{holyoak2018} multiconstraint theory is a natural ally here, not a competitor. Their three constraints on analogy (similarity, structural consistency, and purpose) map onto the account's components: similarity is property overlap in the cluster; structural consistency is relational structure maintained by mechanisms in both domains; purpose is projectibility relative to the reasoner's goals. What the account adds is a causal grounding for why the purpose constraint matters: without goal-directed selectivity, structural overlap is detected but not filtered for projectibility (see Prediction~3 below).

Boyd explains \emph{why} metaphors succeed: accommodation to causal structure. Lakoff explains \emph{how} they arise: embodied experience. The reframing provides what each lacks. The profile account supplies the support that Lakoff's account needs: what maintains property clusters, what makes some metaphors productive (projectible to new cases) and others not, why metaphorical inferences are reliable when they are. Lakoff's embodiment extends Boyd's framework beyond scientific metaphor to the everyday and literary cases Boyd never considered.

Three predictions follow from the reframing, each testable and each diverging from existing accounts:

\begin{enumerate}
\item \textbf{Conventionalization requires mechanism support, not just frequency.} \textcite{bowdle2005} treat the comparison-to-categorization shift as driven by frequency of use. The profile account predicts that frequency is necessary but not sufficient: only metaphors whose relational structure is maintained by mechanisms in the target domain will conventionalize. A frequently used metaphor mapping onto a target with no mechanism support should remain effortful or die. This could be tested by identifying high-frequency metaphors that resist conventionalization and checking whether their targets lack the relevant mechanism support.

\item \textbf{Any animal with profile-structured categories should show cross-domain recruitment.} If metaphor is cross-domain profile extension, the basic operation should be available to non-human animals whose categories track projectibility profiles. Section~\ref{sec:dogs} develops this prediction for dogs.

\item \textbf{Systems detecting structural overlap without goal-directed projectibility constraints should extend indiscriminately.} If cross-domain extension is triggered by structural overlap between property clusters, then a system that detects such overlap but has no goals constraining which extensions matter should project productively and spuriously at equal rates. Large language models fit this description. Their \enquote{spurious connections} are typically genuine structural overlap that isn't projectible for any purpose. The overlap is real; the projectibility isn't. Projectibility is relative to a mechanism suite and a purpose: a geologist and an archaeologist project different property clusters from the same riverbank because they have different investigative goals. Organisms project selectively because they have goals. Systems without goals project everything. They will follow any line~-- even the painted ones. A nuance: when explicitly prompted with a goal (\enquote{produce a novel literary metaphor}), LLMs can be strikingly selective, avoiding conventional source domains and generating structurally novel mappings. The prediction is about default behaviour, not prompted behaviour. Give an LLM a goal and it approximates selectivity; remove the goal and it extends indiscriminately.
\end{enumerate}

These predictions require that \enquote{mechanism support} be assessable, not merely intuited. In practice, mechanism support is estimated through converging proxies rather than read directly off the world: cross-cultural stability of the metaphorical mapping (maintained across unrelated language communities?), robustness of inferential transfer (do the projected inferences hold under variation in context?), default versus effortful processing (is the mapping conventional or does it require active construction?), perturbation response (does the mapping degrade gracefully when source-domain knowledge is disrupted, or collapse entirely?), and diachronic conventionalization trajectory (did the mapping stabilize over time or remain novel?). No single proxy is decisive. The assessment is relative to the property cluster under analysis and the explanatory goals of the investigation, which is a feature of the framework, not a defect: Boyd's own account of natural kinds treats kind-individuation as project-relative \citep[532]{boyd1993metaphor}.

\section{The full range: literary, everyday, constitutive}
\label{sec:literary}

Both Boyd and Lakoff pulled metaphor out of the literary tradition and into cognitive and scientific territory. This paper has followed them. But the framework applies to literary metaphor too, and its prediction there is distinctive.

Literary metaphors typically exploit property overlap with low mechanism support. Dickinson's \enquote{Hope is the thing with feathers} maps hope onto a bird: it alights unbidden, persists through storms, asks nothing in return. The overlap is real but not maintained by shared causal structure. Nothing about the psychology of hope sustains avian behaviour, and nothing about ornithology sustains optimism. \textcite{cooperrider2026} uses the poem as an example of a highly abstract metaphor at the far end of a continuum from basic categorization, and admits he doesn't fully understand it himself. That's the point. Literary metaphors don't need to be fully articulable because they aren't doing inferential work. They're valued precisely because the overlap is unpredicted and aesthetically productive rather than inferentially productive. You don't project from it: hope doesn't migrate south in winter.

The profile account can say something about what makes literary metaphors aesthetically productive, even if it can't say what makes them beautiful. The framework predicts that the best candidates should map between rich property clusters (many internal properties, much relational structure) whose maintaining mechanisms are incommensurable. Dickinson's metaphor works partly because both hope and birds are richly structured kinds, and the mechanisms are utterly different (psychology vs ornithology). The overlap is genuine but can never be stabilized: the reader keeps finding fits (persistence, fragility, unbidden arrival) and misfits (migration, moulting, birdsong), and that tension is the aesthetic effect. A metaphor mapping between impoverished clusters, or between clusters with compatible mechanisms, would lack this productive instability. It would either collapse into a literal comparison or conventionalize into an everyday metaphor. The poet's work, on this account, is not passive discovery of overlap but precise selection at the mechanism boundary: knowing exactly which properties of each cluster to activate and where to stop. The absence of mechanism support is the resource, not the limitation. If mechanisms sustained the mapping, it would conventionalize and stop being interesting. The poet needs the gap, and the art is in choosing exactly the right partial overlap and refusing to extend it one inch further than it can bear.

The framework can push further. In a richly structured profile, properties aren't isolated: they're connected to each other via the mechanisms that hold the cluster together. When Dickinson maps persistence through storms, the reader's knowledge of the source cluster generates the surrounding avian network (migration, moulting, metabolic adaptation) as informative absences. Hope does not, after all, migrate south. Hope does not moult. Each absence illuminates the target by contrast. A merely odd metaphor maps properties with few such connections in either cluster, so the unmapped properties generate no informative contrast: they're just irrelevant. The source cluster's internal connectivity does aesthetic work even where the mapping stops. This is a structural prediction: literary metaphors will be aesthetically productive to the extent that the mapped properties sit at the centre of richly connected clusters, so that the unmapped properties arrive as illuminating contrasts rather than inert remainders. What the framework still can't explain is why one particular pattern of informative absence affects us the way it does. That requires aesthetics, not philosophy of science, and it's the right boundary for the framework.

The grade of support sustaining projectible inference across domains predicts a three-way typology of metaphor. Literary metaphors have the weakest support. They exploit genuine overlap that nothing in the target sustains; they don't conventionalize and aren't meant to. Everyday metaphors (Lakoff and Johnson's domain) are stabilized: embodied experience and shared social practice hold the relevant property clusters in place in both source and target, which is why they conventionalize across communities. \metaphor{argument is war} persists because competition, strategy, and alliance are stabilized in both domains by discursive norms and institutions, not by corrective feedback. Theory-constitutive metaphors (Boyd's domain) sit at the highest grade. The target domain turns out to have its own mechanisms actively maintaining the mapped structure, sometimes controlling it, so thoroughly that the metaphor becomes indispensable to the field and generates a research programme. Strict homeostatic control is the exception in this series, not the rule: the dog's \term{trail} (Section~\ref{sec:dogs}) reaches it; most metaphors that work do so on the stabilized or maintained grades.

The career of metaphor \citep{bowdle2005} maps onto this gradient. A novel literary metaphor that turns out to have mechanism support conventionalizes: it enters the everyday repertoire because the overlap is sustained, not because it's repeated. One that lacks mechanism support stays novel or dies, no matter how frequently it's used. And an everyday metaphor that proves to have deep causal structure can become constitutive: it stops being optional and starts being the only way to ask certain questions. Mechanism support, not frequency alone, drives the trajectory.

\section{What dogs know about trails}
\label{sec:dogs}

The best way to see what this reframing buys is a case study from outside language entirely.

\subsection{\term{Trail} as a projectibility profile}

Dogs possess a functional spatial category that I'll call \term{trail}, distinct from \posscite{talmy2000} and \posscite{lakoff1987} abstract linguistic schema \term{path}. A trail is a physical feature of the landscape. The following properties typically cluster in trails:

\begin{enumerate}
\item Linear extent
\item Visual contrast with surrounding ground
\item Hard or smooth substrate
\item Concentrated conspecific scent
\item Frequent human or animal locomotion
\item Connectivity between locations (source-trajectory-goal structure)
\item Bounded width
\end{enumerate}

These properties co-occur reliably, and the co-occurrence is actively maintained by feedback, which is what makes the cluster homeostatic rather than coincidental. Foot traffic compacts substrate, increasing visual contrast, which makes the trail easier to follow, which concentrates more traffic. Repeated traversal deposits scent, which decays temporally to form a directional sequence, which makes the trail olfactorily trackable, which attracts more traversal. Connectivity between locations provides the functional reason for using the trail, which drives the traffic that maintains every other property. The cluster is self-reinforcing: each property's presence promotes the others. And the perturbation response is diagnostic. If vegetation overgrows a trail segment, continued use clears it. If use drops, all properties degrade together: substrate softens, scent fades, visual contrast disappears. Abandoned trails lose properties at different rates (scent fades first, then visual contrast as vegetation encroaches, while substrate compaction persists longest), but the trajectory is toward dissolution of the full cluster because the maintaining mechanism (use) has ceased. This is homeostasis in the strict, controller sense (Section~\ref{sec:hpc-tcm}): corrective feedback that restores the cluster under perturbation, not mere stabilization. Mechanisms \enquote{favour the presence of some of the properties given the presence of the others} \citep[526]{boyd1993metaphor}, and the perturbation response (overgrowth cleared by use) is what marks the top grade of support. The trail is the clean case where a metaphor-relevant category is held together by genuine control.

Four independent lines of evidence confirm that dogs represent trails with precisely this profile structure.

\textcite{thesen1993} showed that German shepherds exhibit three distinct behavioural phases when tracking: searching, deciding direction, and tracking. During the deciding phase, dogs slow to half speed, examine 2--5 footprints, and triple their sniffing duration while maintaining a constant 6~Hz sniff frequency with the nose held about 1~cm from the ground \citep[250--251]{thesen1993}. This is active structure extraction, not passive chemotaxis. The authors note that \enquote{the deciding phase seems to be the most difficult one; it is certainly the most impressive one from a human point of view} \citep[251]{thesen1993}. The dogs achieved 100\% success across both grass and concrete surfaces and across 3-minute and 20-minute-old tracks, confirming projectibility across substrate and temporal variation.

\textcite{hepper2005} refined the picture by determining that dogs require exactly five sequential footsteps (spanning approximately 1--2 seconds of temporal decay) to determine trail direction. With five footsteps, all six dogs in the study were successful (9--10/10 correct; $p < 0.01$). With three footsteps, all six performed at chance (3--6/10 correct) \citep[296--297]{hepper2005}. The critical finding is that when footstep order was randomized (same total odour, same spatial extent), \enquote{dogs performed at the level of chance\ldots\ They back-tracked much more frequently and often changed the direction in which they tracked} \citep[294]{hepper2005}. And when individual odour was removed but ground disturbance remained (rubber overshoes), performance again dropped to chance \citep[295]{hepper2005}. Sequential order and individual odour are both constitutive, not incidental.

\textcite{cafazzo2012} studied free-ranging domestic dogs in Rome and found that males concentrate scent-marking at crossroads along territorial boundaries ($t = 5.00$, $z = 2.98$, $n = 16$, $p = 0.003$). Over 282.53 hours of observation and 782 urination events, dogs marked territorial boundaries at higher frequency than interior zones ($z = 3.51$, $p = 0.0004$) \citep[960]{cafazzo2012}. The alpha male deposited 81.82\% of defecations at boundaries and crossroads, 77.78\% on conspicuous or elevated substrates. Crossroads are topological nodes where trails intersect. The preferential marking at these sites shows that dogs don't just use trails~-- they recognize and exploit the functional topology of the trail network.

Crucially, dogs don't just track trails~-- they form expectations about what a trail connects to. \textcite{brauer2018} presented dogs with a violation-of-expectation paradigm: a dog followed the scent trail of one toy (Target~A) but found a different toy (Target~B) at the trail's end. On the first trial, dogs showed measurable surprise~-- additional searching for Target~A~-- interpreted as evidence that dogs represent the identity of what they're tracking, not just a gradient to follow \citep[194--195]{brauer2018}. \textcite{brauer2021} extended this to people: dogs followed the scent trail of one familiar person but encountered a different familiar person, and showed significantly greater behavioural excitement in the mismatch condition ($p = 0.015$). If a trail were merely a stimulus gradient, there would be nothing to violate. The violation-of-expectation result shows that dogs represent the trail as connecting to a specific entity at its terminus~-- exactly the kind of structured expectation that a functional category supports.

These aren't independent cues assembled at perception. They're a causally maintained cluster. Walking deposits odour that decays temporally while forming a spatial sequence. The convergence across studies is striking: Thesen found dogs examine 2--5 footprints during the deciding phase; Hepper and Wells found five sequential footsteps required for reliable direction determination (three were insufficient). Both studies independently converge on a small number of discrete odour samples as the unit of directional inference. This is comparative analysis across a structured temporal sequence, not moment-to-moment gradient following.

The claim that \term{trail} is a functional category rather than a product of generic perceptual generalization receives independent support from the spatial navigation literature. \textcite{chan2012} propose a four-part functional taxonomy of landmarks (beacon, orientation cue, associative cue, reference frame), each grounded in distinct neural substrates: beacons in the striatum, orientation cues in the retrosplenial cortex, associative cues in the striatum and parahippocampal gyrus, and reference frames in the hippocampus. The key point: the same visual object can serve different navigational functions depending on context (a church spire functions as a beacon when the church is the goal, an orientation cue when used for heading direction). Function is not read off the percept; it is assigned by the navigational system. This is architectural categorization, not deliberative classification. Boundary information in particular is encoded \enquote{incidentally, even in task manipulations in which encoding of boundary information was neither encouraged nor required} \citep[7]{chan2012}, processed by a dedicated hippocampal system (the \enquote{geometric module} of \citealt{cheng1986}). If dogs have analogous boundary-sensitive spatial processing (and rodent evidence strongly suggests they do), then a painted line would activate boundary/reference-frame processing automatically, not through generic contrast detection.

\subsection{The continuum from trail to \enquote{career path}}

If \term{trail} is a projectibility profile, there should be a continuum from clear members to marginal ones to metaphorical extensions, depending on how many properties are instantiated.

\begin{table}[ht]
\centering\small
\begin{tabular}{@{}>{\raggedright\arraybackslash}p{.24\textwidth}>{\raggedright\arraybackslash}p{.40\textwidth}>{\raggedright\arraybackslash}p{.24\textwidth}@{}}
\toprule
Instance & Properties present & Status \\
\midrule
Well-worn trail through woods & Nearly all (scent, substrate, visual, traffic, connectivity) & Core \term{trail} \\[4pt]
Paved sidewalk beside soccer field & Most (visual, substrate, traffic, connectivity) & Clear \term{trail} \\[4pt]
Painted line on a field & Visual-geometric only (linear, bounded, high-contrast) & Marginal / metaphorical \\[4pt]
\enquote{Career path} & Abstract source-trajectory-goal only & Fully abstract metaphor (humans only) \\
\bottomrule
\end{tabular}
\caption{The continuum from core \term{trail} to metaphorical extension.}
\label{tab:trail-continuum}
\end{table}

Return to the painted line from the introduction. This is the theoretically important case, because it strips away nearly every deflationary alternative simultaneously. There's no scent concentration on painted lines. There's no substrate difference (paint on grass). There's no energy advantage. There's no human to follow. The painted line shares only the visual-geometric properties of a trail: a linear, bounded, high-contrast strip extending through space.

The painted-line observation is naturalistic and unsystematic: my dog, on multiple occasions, tracking along the white lines of a soccer field. No published study has tested whether untrained dogs preferentially follow painted ground lines under controlled conditions. Several confounds would need to be ruled out in any such test: paint odour (field-marking paints are typically latex emulsions with volatile organic compounds that produce a sustained chemical signature, not merely a residual trace), prior reinforcement from leash walks along curbs, subtle handler cues, and the possibility that a high-contrast line functions as a visual boundary to be avoided rather than a path to be followed. Dogs are dichromats with lower spatial acuity than humans, but a 5-centimetre white line on green grass provides ample luminance contrast for detection. \textcite{pelgrim2025}, using head-mounted eye-tracking on dogs during outdoor walks, found that a bottom-up saliency model predicted dog gaze at barely above chance (AUC 0.66); dogs' visual attention was category-driven, not contrast-driven, with pavement fixated 38\% of the time it was in view despite not being the most visually salient surface. This suggests that if dogs attend to painted lines, they are unlikely to be doing so on the basis of low-level salience alone.

The remaining alternative is general edge-following or contrast-tracking behaviour (comparable to thigmotaxis in rodents). Edge-following typically involves locomotion alongside a boundary, not on it. Whether my dog tracks on the line or alongside it is exactly the kind of distinction that requires video analysis and trajectory measurement to resolve. I present the observation as a hypothesis-generating case, not as confirmed evidence. What the observation motivates is a prediction: if \term{trail} is a projectibility profile, then an object sharing only visual-geometric properties with trails should be recruitable by dogs, even when it belongs to a different functional domain. That prediction is testable. Any adequate test would need to resolve three diagnostic questions: (a) does the painted line attract locomotion beyond baseline (not just attention)?, (b) is the attraction visual-geometric rather than olfactory (ruling out paint odour)?, and (c) does the dog travel on the line or alongside it (distinguishing trail-following from edge-following)? The details of how to motivate naturalistic locomotion, manage handler effects, and control for prior reinforcement history are methodological questions for canine cognition researchers, not for a theorist with one dog and a soccer field. What the framework contributes is the theoretical prediction and the diagnostic contrasts, not the protocol.

If the prediction holds, the painted line belongs to a different functional domain from \term{trail}. It's a spatial demarcation artefact, not a locomotion-supporting ecological feature. My dog would have recruited his \term{trail} category to assimilate an object whose geometric profile fits but whose ecological role doesn't. That cross-domain recruitment on partial property overlap is the operation I'm calling metaphorical.

One objection might be that this requires more cognitive sophistication than dogs possess. \textcite{muller2014} provide evidence that it doesn't. Border Collies learning to solve Piaget's support problem (choosing which board to pull to retrieve food) succeeded by tracking a perceptual cue (the vertical level of the reward), not by understanding the causal support relation. When that cue became unreliable in transfer tasks, dogs rapidly switched to alternative perceptual cues within 12 trials ($p = 0.007$), while their performance on the original cue simultaneously declined \citep[1077]{muller2014}. This is structured perceptual category use: the dogs track which properties in a cluster are reliable predictors and update accordingly, without representing why those properties co-occur. Profile extension doesn't require explanatory understanding of the category's maintaining mechanisms. It requires sensitivity to the property cluster. A dog can extend \term{trail} to a painted line by tracking visual-geometric properties (linear extent, bounded width, contrast) without representing what makes trails trails.

\subsection{Wild canids confirm the category}

Comparative canid data confirm that \term{trail} is a genuine functional category, not a behavioural coincidence. \textcite{barja2005} found that Iberian wolves apply zone-dependent marking rules to the same road: near dens, high-density random scat placement on inconspicuous substrates; away from dens, conspicuous elevated placement at crossroads (67.3\% conspicuous vs 8.3\% near dens; $\chi^2 = 23.2$, $p < 0.001$), with the access trail between den and periphery showing intermediate values. The same animal discriminates trail segments by functional zone and adjusts marking accordingly. \textcite{kohn1999} confirm the pattern for coyotes: \enquote{both sexes defaecate equally on trails} \citep[660]{kohn1999}. \textcite{cafazzo2012} bridge wild and domestic canids: free-ranging dogs in Rome concentrated scent-marking at territorial boundaries and crossroads ($z = 2.98$, $p = 0.003$). Trails and junctions are privileged sites for communicative marking across canid species, which presupposes discriminating them from surrounding space.

Movement ecology confirms the pattern at phylogenetic scale. \textcite{fagan2025} analysed GPS tracks from 1,239 carnivores across 16 canid and 18 felid species on six continents and found that canids had 15\% ($\pm$7 CI) greater density of reused travel routeways within home ranges than felids, rising to 33\% ($\pm$16 CI) in shared landscapes. Routeway reuse is a canid-typical behavioural strategy, not a coincidence of particular species or habitats. Domestic dogs show the same pattern: free-ranging dogs in southern Chile's temperate rainforest selected trails and roads when entering forest habitat \citep{sepulveda2015}; free-roaming owned dogs in Arequipa, Peru used dry water channels as ecological corridors, travelling on average seven times farther ($p = 0.002$) and 1.2 times more directionally ($p = 0.027$) than non-users \citep{raynor2020}. Hunting dogs equipped with GPS collars retrace their outbound tracks as a homing strategy in roughly 60\% of 622 field trials \citep{benediktova2020}, treating their own prior trajectory as a reusable navigational structure.

Canid \term{trail} representation has all the hallmarks of a projectibility profile, and at the demanding controller grade: the feedback loop actively restores the cluster under perturbation. No single property is sufficient: odour alone without sequential structure doesn't support direction determination. Mechanisms maintain the cluster: walking, chemical decay, and functional use hold the properties together. The category is projectible: dogs generalize across surfaces, temporal scales, and contexts (pet versus free-ranging versus wild). And the framework predicts boundary cases: an object sharing only visual-geometric trail properties (linear extent, bounded width, contrast) should be recruitable even when it lacks the rest of the cluster. The painted line is the test case.

\section{Phylogeny and what's distinctively human}
\label{sec:phylogeny}

If metaphor is cross-domain profile extension on partial property overlap, then any animal whose categories track projectibility profiles should be capable of the basic operation. Section~\ref{sec:dogs} established that canine \term{trail} is such a profile, at the controller grade. The framework predicts that dogs should recruit \term{trail} for objects in different functional domains that share a subset of the category's properties. Whether they actually do is an empirical question the painted-line observation motivates but doesn't settle.

The boundary between \enquote{broad literal category} and \enquote{metaphorical extension} is sharp on this account (there's a fact of the matter about whether the extension crosses functional domains) but epistemically inaccessible, and that's a prediction, not a problem. A paved sidewalk is a clear trail: it has most of the relevant properties. A painted line is marginal: only visual-geometric properties remain. A \enquote{career path} is fully abstract: only the schematic source-trajectory-goal structure persists. The profile account predicts this gradient: as fewer of the profile's properties are instantiated, typicality decreases and support thins, until at some point the extension reaches across functional domains. That's the boundary between literal and metaphorical. It's determinate but not determinable in practice.

Other animals may show analogous behaviour. Birds mobbing owl decoys extend their \term{predator} category to non-functional perceptual matches. Corvids extend their \term{tool} category to novel materials based on geometric properties. These are suggestive, but the diagnostic difficulty is real: is the decoy in the same functional domain as a predator (threat) or a different one (artefact)? The key criterion is whether the target object belongs to a different functional domain from the source category. The painted line would pass this test clearly if dogs do recruit \term{trail} for it: spatial demarcation is a different functional domain from locomotion, individuated by different maintaining mechanisms (Section~\ref{sec:reframing}).

Even within canids, the pattern may not be limited to \term{trail}. Dogs' dominance and submission displays instantiate the same verticality--dominance correlation that \textcite{lakoff1980} identify as \metaphor{power is up}. Dominance displays recruit vertical extent: raised hackles (piloerection increases apparent height), standing on tiptoe, chin-over-shoulder placement. Submission displays recruit the opposite pole: crouching, belly-exposing, tail-lowering \citep{schenkel1967}. The bidirectionality~-- what \textcite{darwin1872} called the \term{principle of antithesis}~-- lifts this above a single ritualized signal. An entire spatial dimension is being modulated to do work along a social dimension, and the two are maintained by different mechanism suites: physical height by biomechanics, social rank by interaction history and coalition dynamics. The submission case is the more revealing half: there's no biomechanical reason why making yourself low should signal non-aggression. The dog is deploying lowness~-- a spatial property~-- as a communicative act in the social domain. Ethologists describe these displays as ritualized signals, and that framing is compatible: ritualization is the process by which a within-domain correlation (bigger animals tend to win fights) gets decoupled from its biomechanical origin and stabilized as cross-domain communication. The decoupling \emph{is} the cross-domain move. Where \term{trail} extension to a painted line would be ontogenetic, verticality--dominance is phylogenetic: the same basic operation, stabilized by selection.

What's distinctively human isn't the basic operation. It's the superstructure built on top of it. The distinction between reliably expressed and latent category extension turns not on whether the organism possesses the category but on whether mechanisms in the new domain sustain the property cluster in ways that support projectible inference.

\begin{enumerate}
\item Deliberate exploitation. Humans can notice that they're extending a category across domains, recruit the extension intentionally, and deploy it strategically. This is clearest for fresh metaphors; entrenched ones like \enquote{career path} may be processed as conventional categories without any awareness of the underlying mapping \citep[see][]{bowdle2005}.
\item Recursive stacking. Humans build metaphors on metaphors: \enquote{career path} generates \enquote{roadblock on my career path,} which generates \enquote{detour around the roadblock.} Each layer extends the source domain further into the target domain, inheriting inferential structure from the previous extension. There's no evidence that dogs stack extensions this way.
\item Linguistic communication of the mapping. Humans can communicate that they're extending a category metaphorically. Ostensive communication (showing that you're making a metaphorical move) is itself communicative work that requires shared intentionality and a common ground of category knowledge. A dog following a painted line doesn't know it's making a metaphorical extension and can't communicate that fact to other dogs. (No dog has yet circulated \mention{Metaphors We Sniff By}.)
\end{enumerate}

\noindent The absence of evidence for stacking and metacognitive awareness in dogs is not evidence of absence. These are empirical claims about cognitive limitations, not definitional stipulations.

The contrast sharpens if we include systems that detect structural overlap without either goals or metacognition. Large language models identify genuine cross-domain property overlap constantly, but project every such overlap indiscriminately. The result is a three-way contrast: LLMs extend across domains without projectibility constraints (structural overlap triggers extension regardless of purpose)%
\footnote{For a recent illustration, \textcite{imas2026} report that LLMs assigned repetitive tasks produce outputs resembling class consciousness~-- extending the category of worker grievance on surface property overlap (arbitrary management, repetitive labour) without the causal underpinning (material interests, embodiment) that generates genuine labour politics.}; dogs extend selectively, constrained by ecological goals (navigation makes \term{trail} extensions live, social competition makes verticality--dominance extensions live; nothing makes \metaphor{argument is war} extensions live for a dog); humans extend selectively and layer deliberate exploitation, recursive stacking, and linguistic communication on top. What goals add to bare structural overlap is selectivity. What the human superstructure adds to goal-directed extension is awareness and communicability.

Some readers will object that what a dog does isn't \mention{metaphor} but a precursor to metaphor, or a component operation that metaphor exploits. This is a reasonable terminological preference, and nothing in the argument depends on resisting it. The substantive claim is that the basic operation (cross-domain profile extension on partial property overlap) is the same operation in both dogs and humans. Humans layer deliberate exploitation, recursive stacking, and linguistic communication on top. If one restricts \mention{metaphor} to cases requiring intentionality or awareness, then the dog performs a precursor operation that human metaphor elaborates. If one defines \mention{metaphor} by the category operation itself (cross-domain extension), the dog performs metaphor. The theoretical content is identical either way. What matters is the continuity of the underlying operation, not the label.

Whether canine cross-domain extension and human metaphor share this operation by common descent (homology) or by convergent evolution from different cognitive architectures (homoplasy) is a further question the present data can't resolve. The profile reframing holds either way: if the operation is homologous, the account explains why it persists across lineages; if homoplasious, it explains why it converges (the same ecological structure selects for the same category operation). Resolving the homology question would require comparative data across a broader phylogenetic range than canids and primates. Note that the gradualist narrative implied by the three-way contrast (LLMs detecting overlap without selectivity, dogs extending under ecological constraints, humans adding metacognition) is more natural under homology than under convergence, where the canine and human operations would be independently evolved solutions to the same ecological problem and the three-way gradient an artefact of presentation rather than a phylogenetic sequence.

\posscite{bowdle2005} career of metaphor (the comparison-to-categorization shift discussed in Section~\ref{sec:consequences}) maps onto this phylogenetic picture. The basic operation (cross-domain recruitment) is available to any animal with profile-structured categories. What humans add is the capacity to entrench these extensions into stable abstract categories through repeated deliberate use, and to transmit them linguistically. The career of a metaphor, from novel comparison to conventional category, recapitulates in cultural time what the phylogenetic sequence describes over evolutionary time.

\textcite{noyes2019} provide a complementary piece. Generics designate kinds (categories with \enquote{rich causal structure} and \enquote{inductive potential}) but not necessarily essences \citep[2]{noyes2019}. Kindhood and essentialism are separable. This directly supports the profile picture: projectible categories without underlying essences, maintained by causal mechanisms rather than definitions. Generics are themselves a cultural transmission mechanism for such categories: \enquote{generics are a subtle and ordinary linguistic cue by which the kind assumption can be transmitted within a culture} \citep[5]{noyes2019}. The feedback loop~-- kind assumption generates generic use, which reinforces kind assumption~-- is a homeostatic stabiliser.

\section{Implications and conclusion}
\label{sec:conclusion}

For metaphor theory, the profile account provides what CMT has needed: causal grounding (support, not just embodied schemas), a stability explanation (why metaphors persist across communities), a projectibility account (why metaphorical inferences are reliable), failure conditions (why some metaphors don't work), field-relativity (the same phenomenon may support different metaphors for different investigative purposes), and a unified typology (literary, everyday, and constitutive metaphors as points on a single dimension of support grade). \textcite[532]{boyd1993metaphor} stressed that \enquote{cutting nature at its joints} is context-specific: joints depend on the linguistic community and its practical or theoretical projects. This extends directly to everyday metaphor. Different communities may metaphorize the same target domain differently, and both may be legitimate~-- just as jade denotes a genuine gemological kind even though jadeite and nephrite are chemically distinct \citep[396]{boyd1979}.

For the profile account, metaphor reveals a capacity that hasn't been theorized: cross-domain extension. Profile analysis has focused on within-domain categorization~-- biological species, chemical kinds, perhaps social kinds. Metaphor shows that projectibility profiles can reach across domains when partial property overlap triggers recruitment. The sharp-but-epistemically-inaccessible boundary between literal and metaphorical is itself a prediction the account makes but hadn't been asked to deliver. And the phylogenetic depth of the basic operation (cross-domain recruitment predating language) extends its explanatory reach beyond human cognition.

Lakoff's ICMs are hard to falsify. \textcite{glynn2010} notes that ICMs \enquote{produce results that are not systematically falsifiable.} There's no principled way to say \enquote{this isn't a radial category} in Lakoff's framework. The profile account provides exclusion power. If a category has no identifiable support, it's not a projectible profile. If a proposed metaphorical mapping has no mechanisms maintaining the shared relational structure in each domain, or if the mapping doesn't support projectible inference, it's not a productive metaphor. The failure modes are principled and testable.

The reflexivity deserves a final word. This paper's own use of projectibility-profile theory to reframe CMT is a theory-constitutive metaphor in Boyd's sense. It imports inferential machinery from biology (stabilization grades, homeostatic control, property co-occurrence, maintained boundaries) into the study of metaphor and categorization. The import is catachrestic: there was no literal vocabulary for what conceptual metaphors are doing when they succeed. The metaphor is open-ended: we don't yet know all the ways profile-structured thinking about metaphor will generate new questions. And it does constitutive work: it makes questions askable~-- What grade of support maintains this metaphorical mapping? Why does this metaphor project while that one doesn't?~-- that weren't askable before.

The simplest version of the argument is this: Lakoff tells us how metaphors arise; Boyd tells us why they succeed; the projectibility-profile account unifies both by grounding metaphor in property clusters held together by support of some grade, homeostatic control being the strongest. My dog may or may not be extending his \term{trail} category across functional domains. But the framework that the question generated stands on its own merits. It provides what CMT has needed (support, stability, falsification), it provides what profile analysis has lacked (cross-domain extension), and it generates testable predictions: any animal whose categories track projectibility profiles should, under the right conditions, recruit those categories for objects in different functional domains that share only a subset of the category's properties; and systems detecting structural overlap without goal-directed projectibility constraints should extend indiscriminately, which appears to be what large language models do. If a controlled study confirms that dogs follow painted lines as trails, the framework will have passed its first empirical test. If it doesn't, the theoretical synthesis remains, and we'll need a different illustration. My dog will remain unconcerned either way.

\bigskip
\section*{Acknowledgements}

Claude Code (Anthropic, Opus 4.6) was used throughout the development of this manuscript. The tool generated initial prose drafts for all sections, which the author substantially revised, restructured, and in many cases rewrote. The tool also performed literature searches, produced summaries of source texts that the author verified against originals, and handled LaTeX formatting and bibliography management. All theoretical arguments, analytical claims, the selection and framing of evidence, and final editorial decisions are the author's. No content generated by the tool appears without the author's review and revision.

\subsection*{Disclosure of interest}

The author reports there are no competing interests to declare.

\newpage
\printbibliography

\end{document}
